\section{Conclusioni e sviluppi futuri}
\label{sec:conclusioni}

\subsection{Conclusioni}
\label{sec:conclusioni-sintesi}

Il lavoro ha portato in OpenFOAM-13 un modello a due temperature per flussi
ipersonici reagenti, appoggiato alla libreria \mpp{} per tutta la
termochimica. Il solver risolve le equazioni nelle energie, che si conservano
anche attraverso un urto, e ricava le temperature con \mpp; le funzioni
sorgente sono scritte una sola volta e condivise con programmi
zero-dimensionali indipendenti. Dove i modelli della libreria non coincidevano
con quelli del lavoro di riferimento, le sorgenti sono state riscritte sopra di
essa.

Tutti i casi di verifica zero-dimensionali di Casseau
\textit{et al.}~\cite{casseau2016} sono stati riprodotti. Il solver coincide
con i programmi zero-dimensionali entro lo 0{,}3\% in tutti i casi tranne
l'azoto ad alta temperatura, dove arriva all'1{,}2\%, reazioni comprese: il modello è integrato correttamente.
Rispetto al lavoro di riferimento gli stati finali coincidono entro l'1\%, la
composizione della miscela di azoto entro il 2{,}5\% e quella dell'aria entro
il 15\%; l'unica differenza rilevante è un ritardo della temperatura
vibrazionale nelle miscele con azoto atomico, con ogni probabilità dovuto al
modo in cui \mpp{} media i tempi di rilassamento. Il confronto ha anche
mostrato quanto i risultati dipendano da scelte che il lavoro di riferimento non dichiara,
come l'uso dell'energia elettronica o le costanti di velocità per l'aria, e ha
permesso di ricostruirle. Il modello è verificato sul bagno termico, dove non
c'è flusso: per usarlo in casi reali servono gli sviluppi che seguono.

\subsection{Sviluppi futuri}
\label{sec:sviluppi}

Gli sviluppi proposti riguardano il calcolo su più celle, il trasporto
dell'energia e i flussi viscosi. Il primo è quello decisivo e si appoggia agli
altri due.

\begin{enumerate}[leftmargin=1.8em]
  \item \textbf{Calcolo su una griglia vera.} Finora tutto è stato verificato
    su una sola cella: è questo il passo che rende il modello utilizzabile su
    un flusso vero. Che cosa serve è descritto nella
    sezione~\ref{sec:discussione-multicella}: condizioni al contorno per i due
    campi nuovi, sorgenti calcolate con la densità del passo corrente e i
    termini di trasporto del punto successivo. Per il resto non ci sono
    ostacoli, perché tutti i calcoli del modello termodinamico sono locali
    alla cella.
  \item \textbf{Trasporto dell'energia vibro-elettronica.} È la parte del
    punto precedente che si può sviluppare e provare per conto suo. I flussi
    alle facce che servono per il trasporto sono quelli già calcolati per le
    specie, mentre la conduzione usa la conducibilità vibro-elettronica del
    punto successivo.
  \item \textbf{Viscosità e conducibilità a due temperature.} Servono per i
    flussi viscosi e per il calore che raggiunge la parete. \mpp{} le calcola
    già, e divide la conducibilità fra la parte traslazionale e quella
    vibro-elettronica: basta chiederle cella per cella nel modello
    termodinamico, come oggi le energie, e passarle al solver al posto di
    quelle del modello a una temperatura.
\end{enumerate}
